\section{A matching lower bound for dissipative signed local relaxation}
\label{sec:signed-star-signed-lower}

The preceding theorem is an algorithmic upper bound.  We now define a broad
class containing it and prove that no member can improve its polynomial scale
on the same star.

\paragraph{Signed dissipative class $\mathcal R_{\pm}$.}
A member starts from $p^0=0$ and the exact invariant
$r^t=\gamma_\alpha e_o-Hp^t$.  At time $t$ it may adaptively or randomly
choose a coordinate $u_t$ and relaxation $\omega_t\in(0,2)$, then perform
\begin{equation}
  \delta_t=\omega_t r^t_{u_t},
  \qquad
  p^{t+1}=p^t+\delta_t e_{u_t},
  \qquad
  r^{t+1}=r^t-\delta_tHe_{u_t}.
  \label{eq:signed-star-class-update}
\end{equation}
Residuals and iterates may be signed.  The output is the maintained iterate
$p^T$ itself; no uncharged inverse, dense response, or residual postprocessor
is allowed.  Each operation costs $d_{u_t}$.

The interval $(0,2)$ is not an arbitrary aesthetic restriction.  It is
exactly the interval in which every coordinate step dissipates the natural
quadratic error energy.

\begin{lemma}[Exact energy identity]
\label{lem:signed-star-energy}
Let $e^t:=\pi-p^t$ and define
\begin{equation}
  \Phi(e):=\frac12 e^TD^{-1}He.
  \label{eq:signed-star-energy}
\end{equation}
For a coordinate update $p^+=p+\delta e_u$,
\begin{equation}
  \Phi(e)-\Phi(e^+)
  =\frac{\delta}{d_u}\left(r_u-\frac\delta2\right).
  \label{eq:signed-star-energy-identity}
\end{equation}
Thus \eqref{eq:signed-star-class-update} gives
\begin{equation}
  \Phi(e^t)-\Phi(e^{t+1})
  =\frac{\omega_t(2-\omega_t)}{2}
      \frac{(r^t_{u_t})^2}{d_{u_t}}
  \ge0.
  \label{eq:signed-star-energy-descent}
\end{equation}
On the center-seeded star,
\begin{equation}
  \Phi(e^0)=\frac{\alpha}{2m}.
  \label{eq:signed-star-initial-energy}
\end{equation}
\end{lemma}

\begin{proof}
We have $e^+=e-\delta e_u$,
$(D^{-1}He)_u=r_u/d_u$, and
$(D^{-1}H)_{uu}=1/d_u$.  Expanding the quadratic proves
\eqref{eq:signed-star-energy-identity}; substituting
$\delta=\omega r_u$ proves descent.  Initially $e^0=\pi$ and
$H\pi=\gamma_\alpha e_o$, so
\[
  \Phi(e^0)
  =\frac12\pi^TD^{-1}\gamma_\alpha e_o
  =\frac{\gamma_\alpha\pi_o}{2m}
  =\frac{\alpha}{2m},
\]
using $\gamma_\alpha\pi_o=\alpha$.
\end{proof}

\subsection{The fast mode caps every center displacement}

Put
\begin{equation}
  \overline H:=D^{-1/2}HD^{1/2}
  =I-c_\alpha D^{-1/2}AD^{-1/2},
  \qquad y:=D^{-1/2}e.
  \label{eq:signed-star-symmetric-operator}
\end{equation}
Then $\overline H$ is symmetric positive definite and
$\Phi(e)=\tfrac12y^T\overline Hy$.  The normalized adjacency of the star has
the two unit eigenvectors
\begin{equation}
  q_+=\frac{(\sqrt m,1,\ldots,1)^T}{\sqrt{2m}},
  \qquad
  q_-=\frac{(\sqrt m,-1,\ldots,-1)^T}{\sqrt{2m}},
  \label{eq:signed-star-two-modes}
\end{equation}
with eigenvalues $+1$ and $-1$.  Therefore $q_+$ is the slow mode of
$\overline H$ with eigenvalue $1-c_\alpha=\gamma_\alpha$, whereas $q_-$ is
the fast mode with eigenvalue
\begin{equation}
  \mu_+:=1+c_\alpha=\frac{2}{1+\alpha}.
  \label{eq:signed-star-fast-eigenvalue}
\end{equation}
The remaining leaf-difference modes vanish at the center and are irrelevant
to the following displacement argument, even if an adaptive policy breaks
leaf symmetry.

\begin{lemma}[Pathwise center-displacement cap]
\label{lem:signed-star-displacement-cap}
Every center relaxation of every $\mathcal R_{\pm}$ member satisfies
\begin{equation}
  |\delta_t|\le2\sqrt{\alpha(1+\alpha)}.
  \label{eq:signed-star-displacement-cap}
\end{equation}
\end{lemma}

\begin{proof}
Expand $y$ in an orthonormal eigenbasis of $\overline H$.  Energy descent
implies, before and after every operation,
\begin{equation}
  |\langle y,q_-\rangle|
  \le\sqrt{\frac{2\Phi(e^0)}{\mu_+}}.
  \label{eq:signed-star-fast-coordinate-cap}
\end{equation}
A center update changes $e$ by $-\delta e_o$ and hence changes $y$ by
$-(\delta/\sqrt m)e_o$.  Since the center coordinate of $q_-$ is
$1/\sqrt2$,
\begin{equation}
  |\Delta\langle y,q_-\rangle|
  =\frac{|\delta|}{\sqrt{2m}}.
  \label{eq:signed-star-equal-mode-movement}
\end{equation}
The value of $\langle y,q_-\rangle$ lies in the interval
\eqref{eq:signed-star-fast-coordinate-cap} both before and after the update,
so its displacement is at most twice the interval radius.  Therefore
\begin{align*}
  |\delta|
  &\le\sqrt{2m}\,2
       \sqrt{\frac{2\Phi(e^0)}{\mu_+}} \\
  &=2\sqrt{\alpha(1+\alpha)},
\end{align*}
where \eqref{eq:signed-star-initial-energy} and
\eqref{eq:signed-star-fast-eigenvalue} were used in the last step.
\end{proof}

This cap explains the $1/\sqrt\alpha$ scale mechanically.  Moving the center
changes the slow mode that carries the desired PPR mass, but it necessarily
moves the fast mode by the same magnitude.  Dissipativity prevents the fast
mode from storing enough energy to permit a larger center jump.

\begin{theorem}[Signed one-hop class lower bound]
\label{thm:signed-star-signed-lower}
Under \eqref{eq:signed-star-size}, every successful path of every
$\mathcal R_{\pm}$ member performs at least
\begin{equation}
  N_o
  \ge\frac{\pi_o-\varepsilon_{\rm ppr}m}
           {2\sqrt{\alpha(1+\alpha)}}
  \ge\frac{3}{16\sqrt{\alpha(1+\alpha)}}
  \label{eq:signed-star-center-count-lower}
\end{equation}
center relaxations.  Hence its charged work satisfies
\begin{equation}
  W
  \ge mN_o
  \ge\frac{3}
      {256\,\varepsilon_{\rm ppr}\sqrt{\alpha(1+\alpha)}}
  =\Omega\!\left(
      \frac{1}{\sqrt\alpha\,\varepsilon_{\rm ppr}}
    \right).
  \label{eq:signed-star-signed-work-lower}
\end{equation}
The statement is pathwise and therefore also applies to randomized policies
on every path on which they promise success.
\end{theorem}

\begin{proof}
Only a center relaxation changes $p_o$.  Starting from $p_o^0=0$, semantic
accuracy and \eqref{eq:signed-star-center-required} require total signed
center displacement at least $3/8$.  Total absolute center travel is at least
the signed displacement.  Divide by the per-operation cap
\eqref{eq:signed-star-displacement-cap} to obtain
\eqref{eq:signed-star-center-count-lower}.  Each such operation scans the
degree-$m$ center and costs $m\ge1/(16\varepsilon_{\rm ppr})$.
\end{proof}

Combining \Cref{thm:signed-star-semantic-upper,thm:signed-star-signed-lower}
gives the promised-family class theorem
\begin{equation}
  \boxed{
  W^{\star}_{\mathcal R_{\pm}}
  =\Theta\!\left(
      \frac{1}{\sqrt\alpha\,\varepsilon_{\rm ppr}}
    \right)}
  \label{eq:signed-star-class-theta}
\end{equation}
for $0<\alpha<1$ and
$0<\varepsilon_{\rm ppr}\le1/16$, under the exact primitive, output, and
work contract above.  This is the precise sense in which the signed star rung
is promoted to a proved accelerated algorithmic result.
